% !TEX options=--shell-escape
\documentclass[10pt]{article}
\usepackage{amsmath,textcomp,amssymb,geometry,graphicx,enumerate,xcolor}
\usepackage{listings,minted}
\usemintedstyle{perldoc}
\usepackage{caption}
\captionsetup[figure]{labelfont=bf} % Bolded "Figure X:"
\usepackage[font={footnotesize}]{caption} % caption font is smaller
\usepackage{subcaption}
\renewcommand{\thesubfigure}{\Alph{subfigure}}

\captionsetup{
    margin=0.5in,
    justification=justified,  % or raggedright, centering
    font=small,               % optional: adjust caption font size
    labelfont=bf              % optional: makes "Figure 1:" bold
}

% Reduce hyphenation frequency
\hyphenpenalty=5000        % Penalty for hyphenation (default 50)
\tolerance=1000            % Tolerance for line breaking (default 200)
\emergencystretch=3em      % Additional stretch in emergencies

\usepackage{hyperref}
\hypersetup{
    colorlinks,
    citecolor=red,
    filecolor=green,
    linkcolor=blue,
    urlcolor=blue
}
\definecolor{myGreen}{HTML}{008040} % green color

\setlength\parindent{0pt} % no indentation

\graphicspath{
  {"/Users/bfu/Dropbox/Papers/Longitudinal-gen2/Github/gen2-longit/3-figures"}
} % Sets default 

\def\Name{Beverly Fu}  % Your name
% \def\SubjectNo{21}
\newcommand{\SubjectNoPad}{%
  \ifnum\SubjectNo<10 0\SubjectNo\else \SubjectNo\fi
}
\title{Weeklong longitudinal study results for Subject \SubjectNo}
\author{\Name}

% Page setup
\topmargin -0.5in % topmargin = 1in + \topmargin
\textheight 9in
\textwidth 6.5in
\oddsidemargin 0in % left margin = 1in + oddsidemargin
\footskip 1cm % pagenumber to text

\begin{document}
\definecolor{bg}{rgb}{0.95,0.95,0.95}

\maketitle
% \setcounter{section}{-1}

\tableofcontents

\newpage
\section{Study overview}

\subsection{Study setup}

Throughout the study conducted in August 2025, the subject's small intestinal contents were non-invasively sampled 8 times over one week. An evening sampling procedure was used, in which CapScan devices were ingested 3 hours post-dinner, with a paired saliva sample collected at the time of capsule ingestion. In subsequent bowel movements, CapScan devices were retrieved from stool, and a paired stool sample was collected at the time of capsule retrieval.

\begin{figure}[H]
\centering
\includegraphics[width = 6.5in]{Study-design/Study-timeline.pdf}
\caption{\textbf{Study timeline.} A test CapScan device (CapScan0) was ingested on Thursday prior to the continuous weeklong study. A fresh stool sample was provided on Tuesday evening or Wednesday morning/afternoon. A subset of subjects also participated in a concurrent exercise study and ingested two additional CapScan devices 3 hours prior to exercise sessions.}
\label{TransitTime}
\end{figure}

Subjects reported various metadata, including capsule ingestion and retrieval times, food intake, mental health status, and Bristol stool scores.

\subsection{Subject demographics}

\begin{figure}[H]
\centering
\begin{subfigure}[b]{0.15\textwidth}
    \captionsetup{justification=raggedright, singlelinecheck=false}
    \centering
    \caption{}
    \includegraphics[height=4cm]{Study-design/Weeklong-gender-distribution.pdf}
\end{subfigure}
\begin{subfigure}[b]{0.32\textwidth}
    \captionsetup{justification=raggedright, singlelinecheck=false}
    \centering
    \caption{}
    \includegraphics[height=4cm]{Study-design/Weeklong-age-distribution.pdf}
\end{subfigure}
\begin{subfigure}[b]{0.32\textwidth}
    \captionsetup{justification=raggedright, singlelinecheck=false}
    \centering
    \caption{}
    \includegraphics[height=4cm]{Study-design/Weeklong-race-distribution.pdf}
\end{subfigure}
\caption{\textbf{Participant stats.} 20 healthy subjects participated in the study. (A) Sex distribution: 10 female and 10 male. (B) Age distribution: range of 22–46 years, mean 31.3 years. (C) Race: The cohort was predominantly Asian.}
\label{TransitTime}
\end{figure}

\newpage
\section{General data}

\subsection{Sample recovery}

Of the 172 devices administered to participants, 159 were recovered (92\% recovery rate). The vast majority of recovered devices contained sample material (155/159, 97\%), yielding a total of 338 sampled regions.

\begin{figure}[H]
\centering
\includegraphics[width = 5.5in]{Sample-collection/Samples-summary.pdf}
\caption{\textbf{Per subject sample collection.} Depending on device recovery, number of bowel movements, and CapScan sampling efficiency, subjects generated 23–-40 samples. To calculate the average number of regions per device, divide the "Regions" value by the "Devices" value.}
\label{SamplesCollected}
\end{figure}

CapScan devices were cleaned and sectioned into up to 3 regions based on differences in color, viscosity, and air bubble presence. Sample volumes were optimized to enable all downstream analyses 

\begin{figure}[H]
\centering
\begin{subfigure}[t]{0.3\textwidth}
    \captionsetup{justification=raggedright, singlelinecheck=false}
    \centering
    \caption{}
    \label{fig:capsule_classification}
    \includegraphics[width=\textwidth]{Sample-collection/Capsule-classfication.png}
\end{subfigure}
\begin{subfigure}[t]{0.2\textwidth}
    \captionsetup{justification=raggedright, singlelinecheck=false}
    \centering
    \caption{}
    \label{fig:capsule_regions}
    \includegraphics[width=\textwidth]{Sample-collection/Capsule-regions-distribution.pdf}
\end{subfigure}
\caption{\textbf{Regions per device.} (A) Visual classification of capsule contents. \emph{``SI-light"} samples were translucent, often pink or green with food particles. \emph{``SI"} samples were typically golden and not too chunky; a subset appeared red. \emph{``SI-dark"} samples were dark but still liquid. \emph{``Ascending colon (AC)"} samples were sludgy and required squeegeeing from devices. (B) Distribution of regions per device across all samples (mean 2.18 regions per device).}
\label{Capsule_regions}
\end{figure}

\newpage
\subsection{Transit time}

Gastrointestinal transit time is the duration required for ingested material to pass through the digestive system from ingestion to excretion. As each CapScan device is RFID-chipped, transit times were calculated based on each subject's ingestion and collection records. 
\\ \\
Transit time is influenced by numerous factors including diet composition, physical activity, hydration status, gut motility, and individual physiology. Normal small bowel transit time typically ranges from 2--6 hours, while whole gut transit time (mouth to anus) generally ranges from 24–-72 hours in healthy adults.

\begin{figure}[H]
\centering
\includegraphics[width = 2in]{Sample-collection/Transit-time-distribution.pdf}
\caption{\textbf{Summary transit time.} Histogram of transit times in hours of all CapScan devices ingested during the study. Missing devices were not included. Transit times ranged from 12.17--130.97 h (mean 51.9 h = 2.16 days).}
\label{AllTransit}
\end{figure}

\begin{figure}[H]
\centering
\includegraphics[width = 4.8in]{Sample-collection/Transit-time-facet-distribution.pdf}
\caption{\textbf{Summary transit time broken down by individual.} Same data as in Figure \ref{AllTransit}, broken down by subject.}
\label{IndividualTransit}
\end{figure}

\newpage
\subsection{pH}

When possible, the pH of each CapScan region was measured. Surprisingly, the pH of many samples was below the limit of detection of our pH strips. When spot-checked with a more sensitive method, the pH values were very low, suggesting that these capsules sampled gastric contents rather than small intestinal contents.

\begin{figure}[H]
\centering
\includegraphics[height = 4cm]{Sample-collection/pH-distribution.pdf}
\caption{\textbf{pH distribution.} Sample pH values trended toward the lower end.}
\label{pH}
\end{figure} 

\begin{figure}[H]
\centering
\includegraphics[width = 6.5in]{Sample-collection/pH-summary.pdf}
\caption{\textbf{Summary pH.} Heatmap of all CapScan regions by subject, ordered by ingestion order. Samples with pHs presumably lower than 5.5 are marked in red. CapE1 and CapE2 refer to the first and second exercise capsules, respectively.}
\label{pHSummary}
\end{figure}

\newpage
\subsection{Bristol stool scale}
Stool samples were only collected if a CapScan device was located. Self-reported Bristol Stool Scale (BSS) scores were recorded, indicating stool consistency, and range from 1 (hard lumps) to 7 (entirely liquid).

\begin{figure}[H]
\centering
\includegraphics[height = 4cm]{Sample-collection/Bristol-stool-score-distribution.pdf}
\caption{\textbf{Bristol Stool Scale distribution.} Bristol stool scores were mostly around 4.}
\label{BSS}
\end{figure}

\begin{figure}[H]
\centering
\includegraphics[width = 4in]{Sample-collection/Bristol-stool-score-summary.pdf}
\caption{\textbf{Summary Bristol Stool Scale.} Heatmap of all Bristol Stool scores by subject, ordered by CapScan device.}
\label{BSSSummary}
\end{figure}

% \subsection{Cell count}

\newpage
\subsection{Emotional state}
The Positive and Negative Affect Schedule (PANAS) was administered to assess subjects' emotional state. PANAS is a self-report questionnaire consisting of 20 items that measure positive affect (PA) and negative affect (NA) on a scale from 1 (very slightly or not at all) to 5 (extremely). Subjects rated the extent to which they experienced various emotions during a specified time frame.

\begin{figure}[H]
\centering
\includegraphics[width = 6in]{PANAS/PANAS-summary.pdf}
\caption{\textbf{PANAS score summary.} Heatmap of positive affect (interested, excited, strong, enthusiastic, proud, alert, inspired, determined, attentive, active) and negative affect (distressed, upset, guilty, scared, hostile, irritable, ashamed, nervous, jittery, afraid) scores by subject and day.}
\label{PANASSummary}
\end{figure}

\newpage
\subsection{Alpha diversity}

Shannon diversity is a metric that quantifies microbial community diversity by accounting for both the number of species (richness) and their relative abundances (evenness). It is calculated as:

\begin{equation}
H' = -\sum_{i=1}^{S} p_i \ln(p_i)
\end{equation}

where $S$ is the total number of species and $p_i$ is the proportion of the community represented by species $i$. Higher Shannon values indicate greater diversity.
\\ \\
Typical Shannon diversity ranges vary by body site: saliva samples generally range from 2.5--4.0, reflecting moderate to high diversity; small intestinal samples typically range from 1.0-–3.0, with lower diversity due to the harsh digestive environment and rapid transit; stool samples generally range from 3.0–-4.5, representing the highest diversity due to the stable, nutrient-rich colonic environment. These ranges can vary based on sequencing depth, bioinformatic pipelines, and individual factors such as diet, medications, and health status.

\begin{figure}[H]
\centering
\includegraphics[height = 7cm]{Weeklong-Alpha-diversity/Weeklong-Shannon.pdf}
\caption{\textbf{Alpha diversity summary.} Shannon diversity index for all samples, stratified by sample type. Mean values: saliva 3.31, CapScan (small intestine) 3.17, stool 3.55.}
\label{AlphaDivSummary}
\end{figure}


\begin{figure}[H]
\centering
\begin{subfigure}[b]{0.45\textwidth}
    \captionsetup{justification=raggedright, singlelinecheck=false}
    \centering
    \caption{}
    \includegraphics[width=\textwidth]{Weeklong-Alpha-diversity/Weeklong-Shannon-Capscan.pdf}
\end{subfigure}
\begin{subfigure}[b]{0.45\textwidth}
    \captionsetup{justification=raggedright, singlelinecheck=false}
    \centering
    \caption{}
    \includegraphics[width=\textwidth]{Weeklong-Alpha-diversity/Weeklong-Shannon-Capscan-pH.pdf}
\end{subfigure}
\begin{subfigure}[b]{0.45\textwidth}
    \captionsetup{justification=raggedright, singlelinecheck=false}
    \centering
    \caption{}
    \includegraphics[width=\textwidth]{Weeklong-Alpha-diversity/Weeklong-Shannon-Capscan-prediction.pdf}
\end{subfigure}
\begin{subfigure}[b]{0.45\textwidth}
    \captionsetup{justification=raggedright, singlelinecheck=false}
    \centering
    \caption{}
    \includegraphics[width=\textwidth]{Weeklong-Alpha-diversity/Weeklong-Shannon-Capscan-transit-time.pdf}
\end{subfigure}
\caption{\textbf{Alpha diversity of CapScan samples.} Shannon diversity of CapScan samples is influenced by multiple factors. Distribution of Shannon values by (A) capsule region, (B) sample pH, (C) sample appearance classification, and (D) transit time.}
\end{figure}

\begin{figure}[H]
\centering
\includegraphics[width = 6in]{Weeklong-Alpha-diversity/Weeklong-Shannon-subj.pdf}
\caption{\textbf{Alpha diversity broken down by individual.} Same data as in Figure \ref{AlphaDivSummary}, broken down by subject.}
\end{figure}

\newpage
\subsection{Principal coordinate analysis of microbial composition}

Principal Coordinates Analysis (PCoA), also known as classical multidimensional scaling, is an ordination method used to visualize similarities or dissimilarities between microbial communities in low-dimensional space. PCoA takes a distance matrix calculated between samples and represents the samples as points in a coordinate system, where samples with similar microbial compositions are positioned closer together.
\begin{itemize}
  \item \textbf{Binary (Jaccard):} Presence/absence-based metric that considers only whether taxa are present or absent, ignoring abundance information. Emphasizes differences in community membership.

  \item \textbf{Bray--Curtis:} Abundance-based metric that accounts for both the presence and relative abundances of taxa. Ranges from 0 (identical communities) to 1 (completely dissimilar). Does not consider phylogenetic relationships.

  \item \textbf{Canberra:} Abundance-based metric that is sensitive to differences in rare taxa by weighting contributions of each taxon by the sum of their abundances across samples.

  \item \textbf{UniFrac:} Phylogenetic metric that measures the fraction of unique branch length in a phylogenetic tree between two samples. Accounts for evolutionary relationships between taxa. Unweighted UniFrac considers only presence/absence.

  \item \textbf{Weighted UniFrac (wUniFrac):} Phylogenetic metric similar to UniFrac but incorporates abundance information, giving more weight to abundant lineages while still considering evolutionary relationships.
\end{itemize}

\begin{figure}[H]
\centering
\begin{subfigure}[b]{0.32\textwidth}
    \captionsetup{justification=raggedright, singlelinecheck=false}
    \centering
    \caption{}
    \label{fig:pcoa_binary}
    \includegraphics[width=\textwidth]{Weeklong-PCoA/All/PCoA-binary-all.pdf}
\end{subfigure}
\begin{subfigure}[b]{0.32\textwidth}
    \captionsetup{justification=raggedright, singlelinecheck=false}
    \centering
    \caption{}
    \label{fig:pcoa_bray}
    \includegraphics[width=\textwidth]{Weeklong-PCoA/All/PCoA-bray-all.pdf}
\end{subfigure}
\begin{subfigure}[b]{0.32\textwidth}
    \captionsetup{justification=raggedright, singlelinecheck=false}
    \centering
    \caption{}
    \label{fig:pcoa_canberra}
    \includegraphics[width=\textwidth]{Weeklong-PCoA/All/PCoA-canberra-all.pdf}
\end{subfigure}
\begin{subfigure}[b]{0.32\textwidth}
    \captionsetup{justification=raggedright, singlelinecheck=false}
    \centering
    \caption{}
    \label{fig:pcoa_unifrac}
    \includegraphics[width=\textwidth]{Weeklong-PCoA/All/PCoA-unifrac-all.pdf}
\end{subfigure}
\begin{subfigure}[b]{0.43\textwidth}
    \captionsetup{justification=raggedright, singlelinecheck=false}
    \centering
    \caption{}
    \label{fig:pcoa_wunifrac}
    \includegraphics[width=\textwidth]{Weeklong-PCoA/All/PCoA-wunifrac-all.pdf}
\end{subfigure}

\caption{\textbf{Principal coordinates analysis using multiple distance metrics.} PCoA ordination of all samples using (A) Binary (Jaccard), (B) Bray--Curtis, (C) Canberra, (D) UniFrac, and (E) Weighted UniFrac distance metrics. Samples are colored by sample type.}
\label{fig:pcoa_all}
\end{figure}

\newpage
\section{Subject-specific data}
\subsection{Capscan and tubes}

All CapScan devices collected were cleaned and sectioned. The area closest to the white tip is the most proximal region collected while the area closest to the tape-labeled binder clip is the most distal region. We marked targeted sampling regions and took photos of the final samples collected in the tubes. Note that there is generally great variability in terms of the number of regions, colors, and consistencies.

\ifx\pdffilesize{Weeklong-Processed-capsules/20250826_Capsule_Weeklong_S\SubjectNo.png}\undefined
  \vspace{1em} Participant \SubjectNo\ did not collect any capsules.
\else
  \begin{figure}[H]
  \centering
  \includegraphics[width = 5.5in]{Weeklong-Processed-capsules/20250826_Capsule_Weeklong_S\SubjectNo.png}
  \caption{\textbf{Sectioned CapScan devices.} Pictures of cleaned devices and final collection regions.}
  \label{CapscanSamples}
  \end{figure}
\fi

\newpage
\subsection{Transit time timeline}

As each CapScan device is RFID chipped, we were able to calculate transit times. Average transit times per subject can be found in Figure \ref{IndividualTransit}.

\ifx\pdffilesize{Sample-collection/Timeline-breakdown/Timeline-breakdown-\SubjectNoPad.pdf}\undefined
  \vspace{1em} Participant \SubjectNo\ did not provide ingestion and collection times.
\else
  \begin{figure}[H]
  \centering
  \includegraphics[width = 6in]{Sample-collection/Timeline-breakdown/Timeline-breakdown-\SubjectNoPad.pdf}
  \caption{\textbf{Individual CapScan ingestion and collection times.} Timeline of when each CapScan device was in the subject's body throughout the study. The length of each blue bar represents the transit time in hours, which is indicated in white text. If a device was not found in the stool, it is depicted as a transparent bar.}
  \label{Timeline}
  \end{figure}
\fi

\newpage
\subsection{Cell imaging}

During the weeklong study, subjects attempted to bring in a fresh paired capsule and stool sample that was kept refrigerated at 4 \textdegree C. We processed the samples per usual and prepared glycerol stocks and imaged them to look at overall cell densities and cell morphologies.

\ifx\pdffilesize{Weeklong-Cell-count/20250826_Weeklong_Imaging_Subj\SubjectNo.pdf}\undefined
  \vspace{1em} Participant \SubjectNo\ was unable to provide a fresh sample.
\else
  \begin{figure}[H]
  \centering
  \includegraphics[width = 3.7in]{Weeklong-Cell-count/20250826_Weeklong_Imaging_Subj\SubjectNo.pdf}
  \caption{\textbf{Imaging of fresh sample.} Example frames collected for each capsule region and stool sample. Any sample dilution is indicated. Phase contrast is depicted on the left while live/dead staining (indicated by green or red fluorescence, respectively) is depicted on the right.}
  \label{CellImaging}
  \end{figure}
\fi

\newpage
\subsection{Microbial composition}

Microbial community composition was characterized through 16S rRNA gene sequencing performed on the Illumina NextSeq platform. Raw sequencing reads were processed using the DADA2 pipeline, which infers exact amplicon sequence variants (ASVs) by modeling and correcting Illumina sequencing errors. Taxonomic assignment was performed against the SILVA ribosomal RNA database. Samples with fewer than 10,000 reads were excluded from downstream analyses to ensure adequate sequencing depth for reliable community characterization.

\begin{figure}[H]
\centering
\includegraphics[width = 6in]{Weeklong-composition-0p01/Composition-all-legend.pdf}
\caption{\textbf{Taxonomic family legend.} Color key for bacterial families shown in microbial composition figures.}
\label{Legend}
\end{figure}

\begin{figure}[H]
\centering
\begin{subfigure}[b]{\textwidth}
    \captionsetup{justification=raggedright, singlelinecheck=false}
    \centering
    \caption{}
    \includegraphics[height=6.5cm]{Weeklong-composition-0p01/Saliva-collection/Composition-saliva-\SubjectNoPad.pdf}
\end{subfigure}
\begin{subfigure}[b]{\textwidth}
    \captionsetup{justification=raggedright, singlelinecheck=false}
    \centering
    \caption{}
    \includegraphics[height=6.5cm]{Weeklong-composition-0p01/Capscan-collection/Composition-capscan-\SubjectNoPad.pdf}
\end{subfigure}
\begin{subfigure}[b]{\textwidth}
    \captionsetup{justification=raggedright, singlelinecheck=false}
    \centering
    \caption{}
    \includegraphics[height=6.5cm]{Weeklong-composition-0p01/Stool-collection/Composition-stool-\SubjectNoPad.pdf}
\end{subfigure}
\caption{\textbf{Microbial composition.} (A) Saliva sample composition collected at capsule ingestion. (B) CapScan and (C) stool sample composition collected at capsule retrieval. Taxa with relative abundances $>$1\% are shown; all others grouped as "Other." Samples with fewer than 10,000 reads are labeled in red.}
\label{fig:composition}
\end{figure}

\newpage
\subsection{ASVs heatmap}

\begin{figure}[H]
\centering
\includegraphics[width = 8in]{Weeklong-ASV-heatmap-0p01/Saliva-collection/ASV-saliva-\SubjectNoPad.pdf}
\caption{\textbf{ASVs in saliva samples.} Only ASVs detected at $>$1\% relative abundance in at least one sample are shown.}
\end{figure}

\begin{figure}[H]
\centering
\includegraphics[width = 8in]{Weeklong-ASV-heatmap-0p01/Capscan-collection/ASV-capscan-\SubjectNoPad.pdf}
\caption{\textbf{ASVs in CapScan samples.} Only ASVs detected at $>$1\% relative abundance in at least one sample are shown.}
\end{figure}

\begin{figure}[H]
\centering
\includegraphics[width = 8in]{Weeklong-ASV-heatmap-0p01/Stool-collection/ASV-stool-\SubjectNoPad.pdf}
\caption{\textbf{ASVs in stool samples.} Only ASVs detected at $>$1\% relative abundance in at least one sample are shown.}
\end{figure}

\end{document}
